\newpage
\section{Messverstärker}\label{sec:1}

\subsection{Messen der Offsetspannung}\label{sec:1-1}
\subsubsection{Aufgabenstellung}

Zu ermitteln ist die Offsetspannung ($U_{OS}$) des OP-A3 ($\mu$A741 Biploar)
auf dem Messbrett mittels geeigneter Schaltung. Anschließend soll $U_{OS}$
mittels dem zu A3 zugehörigen Potientiometer abgelichen werden.

\subsubsection{Versuchsaufbau}

Zur Messung wird ein nicht invertierender Verstärker als Schaltung verwendet. Alle
Bauteile sind auf dem Messverstärkerbrett (siehe
Geräteliste~\ref{sec:geraeteliste}). Die geziegte Spannungsquelle $U_{OS}$
modelliert die Offsetspannung des OPs. Am Aufbau wird der nicht invertierende
Eingang auf Masse gelegt.

\begin{figure}[h]
	\centering
	\begin{circuitikz}[scale=0.9, transform shape]
			% siehe PR3/latex/main-report.tex 
			\myNIA{A3}{(0,0)}{$R_2$}{$R_1$};
			\draw (A3-in) to[voltage source, v_={$U_{OS}$}] (A3-in |- A3-gnd) node[ground]{};
			\draw (A3-out) to[open, v^={$U_{a,OS}$},o-o] (A3-out |- A3-gnd) node[ground]{};
	\end{circuitikz}
	\caption{Messung Offsetspannung $U_{OS}$}\label{fig:offsetspannung}
\end{figure}

\subsubsection{Berechnungen der Verstärkung}

\begin{gather}
	A=\frac{R_1+R_2}{R_2}\approx \frac{R_1}{R_2}
\end{gather}

Da in der Tabelle~\cite[S.~77, Abb.~3.13]{EMTPraktikum2025} steht,
dass A3, $U_{OS}=\SI{1}{\milli\volt}$ besitzt ist eine Verstärkung von $1000$ bei der Versorgung $\pm\SI{12}{\volt}$ geeignet.

Demnach wurden die Werte für $R_1$ und $R_2$ gewählt:
\begin{gather}
	R_2=\SI{1}{\kilo\ohm} \\
	R_1 = A \cdot R_2 = 1000 \cdot \SI{1}{\kilo\ohm} = \SI{1}{\mega\ohm}
\end{gather}

\subsubsection{Messung}

Bei der Messung der Offsetspannung am Ausgang vom OP, ist es notwendig den
Taster neben den Potentiometers zu betätigen, um das Potentiometer zu
überbrücken. Ansonsten würde dieser $U_{OS}$ Verändern.

Mit dem Multimeter Metex (20V DC Messbereich, siehe Geräteliste~\ref{sec:geraeteliste}) wurde $U_{a,OS}=\SI{3.45}{\volt}$ gemessen.
\begin{gather}
	U_{OS} = \frac{U_{a,OS}}{A} = \frac{\SI{3.45}{\volt}}{1000} = \SI{3.45}{\milli\volt}
\end{gather}

\subsubsection{Abgleichen der Offsetspannung}

Nun wird der Taster nicht mehr gedrückt und weiter mit dem Multimeter
$U_{a,OS}$ gemessen. Das Potentiometer wird so eingestellt, bis beim kleinsten
Messbereich vom Multimeter~\ref{sec:geraeteliste} (200mV) die Spannung 0 wird,
bzw. bis sie nicht mehr kleiner wird. Im Anschluss wurde eine Spannung von $\SI{0.7}{\milli\volt}$ gemessen.

\subsubsection{Erkenntnisse}

Zu sehen ist, dass die Offsetspannung $U_{OS}=\SI{3.45}{\milli\volt}$ nahe an dem im
Skriptum angegebenen typischen Wert von $\SI{1}{\milli\volt}$ liegt. Die Abgleichung
muss für OPVs in den folgenden Versuchen immer durchgeführt werden.

